\documentclass[11pt]{article}
\usepackage[margin=1in]{geometry}

\begin{document}

\begin{center}
{\large Short-Term Temperature Forecasting with Jena Climate Data}\\
CSCI 6840 Final Project Part 2 Proposal
\end{center}

\noindent\textbf{0. Project Idea / Title.}
This project will use data mining methods to predict future temperature from historical weather measurements. The problem is a supervised regression task for time-series forecasting: given previous readings, predict temperature 1 hour, 6 hours, and 24 hours ahead. We will compare temperature-only forecasting with multivariate forecasting to test whether pressure, humidity, wind, dew point, and related variables improve accuracy.

\medskip
\noindent\textbf{1. Team Members.}\\
Kareem Washington, washingtonk21@students.ecu.edu\\
John Aramis Scott, scottjo24@students.ecu.edu

\medskip
\noindent\textbf{2. Data Source.}
We will use the public Jena Climate dataset, which contains about 420,000 readings from 2009-2016. The data are recorded every 10 minutes and include 14 weather variables such as temperature, pressure, humidity, wind speed, wind direction, and dew point. This fits the course topic of time-series data because observations are indexed in time order and nearby readings are related.

\medskip
\noindent\textbf{3. Objectives.}
Our objective is to predict temperature at 1-hour, 6-hour, and 24-hour horizons. We will compare a persistence baseline, a temperature-only model, and multivariate models. Success will be measured with MAE and RMSE, with optional $R^2$. A successful result means the learned models reduce MAE/RMSE compared with the persistence baseline and show whether extra weather variables add useful information beyond past temperature.

\medskip
\noindent\textbf{4. Motivation (Why It Matters).}
Short-term weather prediction supports agriculture, transportation, energy planning, and daily decision making. The course emphasizes turning raw data into useful knowledge. Weather readings become more valuable after cleaning, transformation, visualization, and modeling, and this project can also explain seasonal patterns and relationships among measurements.

\medskip
\noindent\textbf{5. Approach (Initial Plan).}
We will follow the course KDD process: data cleaning, data selection, data transformation, data mining, evaluation, presentation, and applying knowledge. First, we will inspect missing or bad readings, verify time order, and decide whether to resample the 10-minute data to hourly data. Next, we will create lag/window features and explore correlations, seasonal patterns, histograms, and scatter plots. For modeling, we will start with a persistence forecast, then train linear regression and random forest regression models using Python, pandas, NumPy, scikit-learn, and Matplotlib. We will use chronological training, validation, and test splits to avoid data leakage. If time allows, we may add anomaly detection for unusual weather patterns or a neural-network model for multi-step forecasting.

\end{document}
